\section{AI for Math 的评估路线：从题库到研究贡献}

如果要判断 AI for Math 是否真正进步，不能只看模型做了多少竞赛题。竞赛题有价值，因为答案清晰、反馈硬；但它只是入口。更高层次的评估应该看形式化证明通过率、定理库贡献、lemma 发现、反例构造、数学家协作效率和新研究方向的产生。

\subsection{评估层级}

\noindent\begin{longtable}{p{0.22\textwidth}p{0.35\textwidth}p{0.34\textwidth}}
\toprule
层级 & 可测指标 & 局限 \\
\midrule
竞赛题求解 & 正确率、pass@k、解题时间 & 容易过拟合题库，不代表研究能力。 \\
形式化补全 & Lean theorem 通过率、tactic 成功率 & 依赖定理库覆盖和题目选择。 \\
证明生成 & 完整证明长度、可读性、验证通过率 & 正确不等于优雅，长证明可能没有洞察。 \\
定理库贡献 & 新 lemma 数量、复用率、维护成本 & 需要社区接受，不只是机器生成。 \\
研究协作 & 数学家节省时间、发现反例、提出猜想 & 难以标准化，但最接近真实价值。 \\
\bottomrule
\end{longtable}

\begin{importantbox}{评估路线的核心}
AI for Math 的评估应从“答案正确”逐步走向“证明可验证”“证明有解释力”“定理库可复用”“数学家愿意协作”。越往后越难量化，但也越接近真正价值。
\end{importantbox}

\subsection{为什么这和 Agent 相关}

AI for Math 看似是科学问题，实际上也是 Agent 问题：模型要进入 Lean 环境，读取 goal state，调用 tactic，观察错误，修复证明，再继续推进。这和 web agent、coding agent 的闭环非常相似，只是环境更形式化、反馈更硬。EP137 因此可以看作 EP139 Agent 技术史的一种科学版实例。

\begin{knowledgebox}{Lean 环境就是一种高质量 Agent 环境}
Lean 提供明确状态、明确动作、明确验证和明确错误反馈。相比开放网页环境，它更严格；相比普通聊天，它更可训练。这使 AI for Math 成为研究 Agent 学习能力的理想试验场。
\end{knowledgebox}

\subsection{本章小结}

AI for Math 的评估不能停在题库正确率。真正的路线，是从解题、形式化、证明生成，逐步走向定理库贡献和数学家协作。


